\chapter*{Preface}
\addcontentsline{toc}{chapter}{Preface}

These notes accompany the graduate course \emph{Agentic Development in Research}.
The course trains researchers to build, orchestrate, and evolve agentic AI systems for use in real research work.
These notes collect the conceptual foundations that lectures and studio sessions draw on.
They do not replace hands-on practice with the tools; they equip students to understand what they are doing and why.

The notes follow a \emph{capability ladder}: students progress from \emph{operating} a professional agentic rig, to \emph{building} a research agent, to \emph{orchestrating} multi-step workflows, and finally to \emph{evolving} artifacts through automated search.
Two conceptual threads run the length of this ladder.
The first is \emph{evaluation}: the course treats automatic, honest evaluation as the foundation of every trustworthy agentic system, and reveals how an evaluation harness becomes a fitness function once the system evolves.
The second is \emph{context engineering}: the discipline of deciding what information the agent holds in its context window, and constructing that context deliberately.

The expected background is some Python experience and general quantitative ability.
No prior knowledge of agents, the shell, or professional development tooling is assumed; the course starts there.
Examples are drawn from \emph{model discovery from data} (symbolic regression), a domain-portable task that researchers across quantitative fields can connect to their own work.

A note on formalism: where these notes adopt mathematical notation, the goal is sharpness of intuition rather than formal rigor.
A definition that names something precisely is more useful than three paragraphs of informal description, even when the thing being named is fundamentally a system to be built rather than a theorem to be proved.
